\section{The term structure: curves, forwards, returns and curve risk}

These entries build spot, par and forward curves from bond prices, connect forwards to expected and realised returns, describe the swap curve and short-term spreads, set out the traditional theories, and measure exposure to curve factors. Continuous-compounding forward rates and FRAs are in \hyperref[mod:forward-rates-fra]{forward rates and FRAs}, term-premium estimates in the same entry, and carry and rolldown trades in \hyperref[meth:curve-trades-carry-rolldown]{curve trades}. Numbers are checked in \texttt{notes/artefacts/term-structure-dynamics/term-structure-checks.py}; annual compounding unless stated.

\begin{qfmethod}[Spot, par and forward rates: definitions, bootstrapping spot rates from par yields, pricing a bond with spot or forward rates, converting between the three curves, and comparing their shapes]\label{meth:spot-par-forward-curves}
\qffield{Statement} The \emph{spot} (zero) rate $s_n$ discounts a single cash flow at $n$: $P_n=100/(1+s_n)^n$. The \emph{par} rate for maturity $n$ is the coupon that prices a bond at par; the \emph{forward} rate $f_{m,n}$ is the rate for a loan from $m$ to $n$ agreed today, $(1+f_{m,n})^{n-m}=(1+s_n)^n/(1+s_m)^m$, and can be locked in by buying the $n$-year zero and shorting the $m$-year zero in equal values \citep[chap.~31, pp.~692--693]{s7-fabozzifixed}. Spot rates come from coupon bond prices by \emph{bootstrapping}: each bond is a package of zeros, so with the shorter spot rates known, the price of the next bond leaves one unknown (chap.~3, pp.~45--49; \citep[sec.~4.7, pp.~106--108]{s7-hulloptions}). A bond priced with spot rates, or with the chain of forward rates, gives the same value: $(1+s_n)^n=\prod_{k=1}^n(1+f_{k-1,k})$. Arbitrage by stripping and reconstitution keeps traded bonds close to their spot-rate values (chap.~3, pp.~49--51).

\qffield{Practice} Example, original. Annual-pay par yields of 3.80, 4.00, 4.15, 4.25 and 4.30 percent for one to five years bootstrap to discount factors 0.963391, 0.924485, 0.884929, 0.846193 and 0.809572, so:
\begin{itemize}
\item spot rates 3.8000, 4.0040, 4.1591, 4.2636 and 4.3155 percent;
\item one-year forward rates $f_{n-1,n}$ of 3.8000, 4.2084, 4.4700, 4.5777 and 4.5235 percent;
\item the implied spot curve one year ahead, $f_{1,n}$, 4.2084, 4.3391, 4.4186 and 4.4448 percent for two to five years.
\end{itemize}
Every par bond reprices to 100 from the discount factors, and the par yields are recovered as $(1-DF_n)/\sum_{k\le n}DF_k$. A 5-year 6 percent bond is worth 107.5286 by spot rates, by forwards and by discount factors; its yield, 4.2948 percent, lies below the 5-year spot rate because its coupons are discounted at lower short rates. Hull's bootstrap from bills and two coupon bonds gives 1.603, 2.010, 2.225, 2.284 and 2.416 percent (continuous) (pp.~107--108).

\qffield{Conditions} \emph{Comparing the curves.} Spot rates are geometric averages of forwards, and par rates averages of spot rates, so on an upward-sloping curve forward $\ge$ spot $\ge$ par, with the order reversed on an inverted curve; the curves coincide only when flat (chap.~31, pp.~694--695). Forwards magnify changes in slope: in the example the 5-year forward (4.5235 percent) is below the 4-year forward while spot and par still rise. The par curve is the flattest, and on-the-run and swap curves approximate par curves (pp.~692 and~694).

\qffield{Pitfalls} Bootstrapping across bonds with different credit, liquidity or tax treatment, or including on-the-run issues whose special liquidity distorts their yields (chap.~3, p.~45, n.~1). Treating an interpolated spot curve as data: the interpolation between instruments is a modelling choice (Hull, p.~108). Reading a forward rate as a forecast (\hyperref[meth:forward-rates-expected-returns]{forwards and returns}).

\qffield{See also} \hyperref[def:interest-rate-measures]{interest-rate measures}; \hyperref[mod:forward-rates-fra]{forward rates and FRAs}; \hyperref[data:gsw-treasury-curve]{the Fed's fitted curve}.

\end{qfmethod}

\begin{qfmethod}[Forward rates and bond returns: forwards as break-evens, the rolling yield, realised against expected returns, assumptions about how spot rates evolve relative to forwards, and riding the yield curve]\label{meth:forward-rates-expected-returns}
\qffield{Statement} \emph{Break-evens.} The implied spot rates one year forward are the future spot rates at which every bond earns the same one-year return as the one-year zero. If forwards are realised, all government bonds earn the same holding-period return and every self-financed curve position breaks even. If the curve is unchanged, an $n$-year zero earns the forward $f_{n-1,n}$, its \emph{rolling yield}: yield income plus the \emph{rolldown} gain as it ages down an upward-sloping curve \citep[chap.~31, pp.~695 and~711--712]{s7-fabozzifixed}. Approximately, $f_{n-1,n}\approx s_n+(n-1)(s_n-s_{n-1})$ (eq.~(31C.2)). Expected return decomposes into yield income, rolldown, the value of convexity and the duration effect of an expected curve change (chap.~32, pp.~715 and~724--725).

\emph{Assumptions in active management.} A manager who expects forwards to be realised expects no bond to beat cash. One who expects an unchanged curve (rates follow a random walk) expects longer bonds to earn their forwards, and so to beat cash when the curve slopes up. Ilmanen reports that the unchanged-curve forecast has been more accurate than the forwards on average, and that steeper curves have been followed by larger excess returns, against the pure expectations hypothesis (pp.~702--703 and~712).

\qffield{Practice} \emph{Reproduced.} Ilmanen's example: one- and two-year spot rates of 6 and 8.08 percent imply $f_{1,2}=10.20$ percent; if it is realised, the two-year zero rises from 85.61 to 90.74, a 6 percent return, the same as the one-year zero. The one-year rate must rise by more than 420 basis points before the two-year zero underperforms (p.~695). His approximation gives 10.16, 13.00, 14.28 and 13.76 percent against exact forwards of 10.20, 13.07, 14.36 and 13.77 (p.~713).

\emph{Original} (the curve of \hyperref[meth:spot-par-forward-curves]{spot, par and forward rates}). If forwards are realised, the two- to five-year zeros all earn 3.80 percent in the first year; if the curve is unchanged, each $n$-year zero earns $f_{n-1,n}$. Riding the curve, buying the 5-year zero and selling it after a year, returns 4.5235 percent if the curve is unchanged against 3.80 percent on the 1-year zero, and 2.5423 percent if all rates rise 50 basis points. The ride breaks even at an 18.1 basis point rise in the 4-year rate, the forward-implied change $f_{1,5}-s_4$.

\qffield{Conditions} The identities are exact for zero-coupon bonds and default-free cash flows; for coupon bonds the rolling yield and forwards differ slightly because of reinvestment assumptions (chap.~32, p.~724, n.~7). Realised returns depend on the actual path of rates, reinvestment and the sale price (\hyperref[meth:bond-return-decomposition]{sources of return}). Expected returns need a model of term premia and convexity; neither is observable (chap.~31, pp.~701--702).

\qffield{Pitfalls} Taking a bearish position whenever rates are expected to rise: it pays only if they rise by more than the forwards imply, because the long bond's positive carry is the cushion (pp.~703--704). Treating rolldown as risk-free income: it assumes the curve does not change. Calling a forward an expectation of the future spot rate (\hyperref[res:term-structure-theories]{theories}).

\qffield{See also} \hyperref[meth:curve-trades-carry-rolldown]{curve trades, carry and rolldown}; \hyperref[def:bond-yield-spread-measures]{yield to maturity}.

\end{qfmethod}

\begin{qfdefinition}[The swap curve and swap spreads; short-term spreads that gauge economy-wide credit and liquidity risk: TED, LIBOR--OIS and commercial-paper tier spreads]\label{def:swap-curve-short-term-spreads}
\qffield{Statement} \emph{The swap curve.} The fixed rates on par interest-rate swaps for a range of maturities form a par curve; the swap curve built on receive-fixed, pay-floating contracts is ``by construction a par curve'' \citep[chap.~31, p.~692]{s7-fabozzifixed}. Market participants use it because swaps trade in every maturity, in size and with flexible dates, making it ``the fair value baseline curve'' for relative value and a benchmark for pricing and hedging (chap.~35, p.~814). OIS rates out to a year give risk-free zero rates directly; longer OIS rates define par bonds and are bootstrapped like a par curve; with SOFR replacing LIBOR, the SOFR OIS curve has become the U.S. dollar swap curve \citep[sec.~7.2, pp.~175--176]{s7-hulloptions}.

\emph{Swap spread.} The swap spread for a maturity is the swap rate minus the Treasury yield of that maturity. Before 2008 most swap spreads were positive, reflecting the view that receiving fixed under a contract was riskier than owning a Treasury; since the crisis many have been negative for long periods, attributed to heavy Treasury supply, central clearing that reduced swap counterparty risk, and capital rules that reduce dealers' willingness to hold Treasuries against paid-fixed swaps \citep[chap.~64, pp.~1580--1581]{s7-fabozzifixed}.

\emph{Short-term spreads.} The TED spread, three-month Eurodollar (LIBOR) deposit rate minus the three-month Treasury bill rate, was 30 to 50 basis points in normal conditions and exceeded 450 in October~2008; the three-month LIBOR--OIS spread briefly reached 364 basis points that month, showing banks' reluctance to lend beyond overnight \citep[sec.~8.2, p.~209]{s7-hulloptions}. OIS carries little bank credit risk, so LIBOR--OIS isolated bank term credit and liquidity premia \citep[chap.~13, p.~305]{s7-fabozzifixed}.

\qffield{Conditions} With U.S. dollar LIBOR ended (\hyperref[meth:policy-path-ois-futures]{LIBOR cessation}), TED and LIBOR--OIS are historical series. Within commercial paper, the spread between the two top rating tiers of non-financial paper, reported by the Federal Reserve, has a median of about 25 basis points, fell to single digits before the 2007--2009 crisis and spiked above 600 at its depth; commercial paper also trades at a spread to Treasury bills that reflects credit risk, collateral use and state taxation of private interest (pp.~305--306).

\qffield{Practice} Example, original (invented annual SOFR OIS par rates of 3.60, 3.50, 3.55, 3.65 and 3.75 percent). The rates bootstrap to discount factors that price each par swap's fixed leg at par, with one-year forwards of 3.60, 3.40, 3.66, 3.97 and 4.19 percent, a dip priced for policy easing. Against a 4.05 percent 5-year Treasury yield the 5-year swap spread is $3.75-4.05=-30$ basis points. A corporate bond's I-spread is measured against the swap curve (\hyperref[def:bond-yield-spread-measures]{spread measures}).

\qffield{Pitfalls} Reading a negative swap spread as negative bank credit risk; it reflects supply, balance-sheet and clearing effects. Mixing a LIBOR-era swap curve with SOFR-era bond spreads. Reading a swap curve as pure expectations (\hyperref[def:interest-rate-swaps]{interest-rate swaps}).

\qffield{See also} \hyperref[mod:multi-curve-fra-swap-valuation]{multi-curve valuation}; \hyperref[def:money-market-reference-rates]{money-market reference rates}; \hyperref[meth:swap-valuation-replication]{swap valuation}.

\end{qfdefinition}

\begin{qfresult}[Traditional term-structure theories: pure expectations in its local, unbiased and return-to-maturity forms, liquidity preference, preferred habitat and market segmentation; what each implies for forward rates and curve shape, the convexity bias and the evidence]\label{res:term-structure-theories}
\qffield{Statement} \emph{Pure expectations.} Forward rates exclusively represent expected future rates, so a rising curve means rates are expected to rise and a flat one that they are not \citep[chap.~3, p.~55]{s7-fabozzifixed}. Its interpretations are not equivalent: that all bonds earn the same expected return over any horizon; the \emph{local} form, that they earn the same expected return over the next short period, the only form Cox, Ingersoll and Ross found sustainable in equilibrium; and the \emph{return-to-maturity} form, that rolling short bonds earns the same as holding a zero to the horizon (pp.~56--58). The \emph{unbiased} form says forward rates equal expected spot rates; with uncertainty it differs from the local form by the convexity bias (chap.~31, p.~701, n.~13).

\emph{Biased expectations.} The \emph{liquidity theory} adds a premium that rises with maturity, so forwards exceed expected rates and an upward curve is consistent with flat or even falling expected rates. The \emph{preferred-habitat theory} allows a premium of either sign to draw investors out of the maturities their liabilities favour, so any shape is possible. The \emph{market-segmentation theory} holds that participants do not leave their sectors, so each maturity's yield is set by its own supply and demand (chap.~3, pp.~58--59). Hull's account is in \hyperref[mod:forward-rates-fra]{forward rates and FRAs}.

\qffield{Practice} \emph{The convexity bias}, original. The one-year rate is 5 percent; next year it will be 3 or 7 percent with equal probability. If the local expectations hypothesis holds, the two-year zero earns 5 percent expected over the first year and the forward rate is 4.9619 percent, 3.8 basis points below the 5 percent expected rate. If instead the forward equalled the expected rate (unbiased expectations), the two-year zero would expect 5.0381 percent, more than the one-year rate. Ilmanen's illustration with 100 basis point annual volatility shows the bias growing with convexity, so long forwards slope down even when rates are expected flat (chap.~31, pp.~699--701). \emph{A liquidity premium}: flat expected rates of 4 percent plus a premium rising 0.10 point per year of maturity give an upward-sloping spot curve.

\qffield{Conditions} \emph{Evidence.} Upward-sloping curves have been followed on average by small declines in long rates, not the rises pure expectations predicts, and steeper curves by larger excess returns; constant-premium versions also fail, because premia vary over time and can be negative (pp.~702--703). Model estimates of the term premium for 10-year Treasuries (\hyperref[mod:forward-rates-fra]{ACM and Kim--Wright}) are on average positive but negative in many months since 2015. Those estimates depend on the model.

\qffield{Pitfalls} Inferring the market's expected policy path from forwards without subtracting a risk premium; at short maturities the premium biases the inference upward (p.~702). Treating the theories as competing facts; each is a restriction that the data reject in its strict form.

\qffield{See also} \hyperref[meth:forward-rates-expected-returns]{forwards and returns}; \hyperref[meth:policy-path-ois-futures]{policy path from OIS and futures}.

\end{qfresult}

\begin{qfmethod}[Yield-curve factor exposures: level, slope and curvature from principal components, factor and key-rate durations, hedging with two instruments, the maturity structure of yield volatility, and economic drivers of views on rates, spreads and curve shape]\label{meth:yield-curve-factor-exposures}
\qffield{Statement} \emph{Factors.} Principal-components analysis of yield changes finds a few independent shifts that explain most of their variance: a level shift (all yields one way), a slope shift (short and long yields opposite) and a curvature or hump shift; for daily U.S. Treasury yields of one to 30 years in 1984--2020 the first explained about 92 percent, the second 5 and the third 1 \citep[chap.~33, pp.~745--747]{s7-fabozzifixed}. The first two are stable across periods and countries; the third and fourth vary with the maturities used and over time, and near the zero bound the first component became upward-sloping because the short end could not move (pp.~748--757). \emph{Exposure} to a factor is the sum over maturities of a position's sensitivity to each rate (its key-rate durations, \hyperref[meth:curve-based-empirical-duration]{curve-based measures}) times the factor's loading; slope duration is derived the same way (pp.~770--773).

\emph{Economic drivers.} In a standard macro model, level shifts come from changes in expected long-run inflation or real rates, slope shifts from cyclical shocks to the output gap that move the short rate, and curvature from changes in the speed at which rates are expected to adjust; the term premium and the central bank's policy bias, both time-varying, also move the curve (pp.~760--767). Ilmanen adds expectations, bond risk premia and the convexity bias as the three influences on curve shape (chap.~31, pp.~695--702). Credit spreads respond to the level and slope of the Treasury curve, the business cycle and confidence as well as to issuer news (chap.~10, p.~254).

\qffield{Practice} \emph{Reproduced} \citep[sec.~22.9, pp.~534--537]{s7-hulloptions}. Hull's daily Treasury factors (2010--2020) have score standard deviations whose squares give 87.3 percent of variance to the first factor, 95.6 percent to the first two and a further 2.1 percent to the third. A portfolio changing by +10, +4, $-8$, $-7$ and +2 million dollars per basis point at 2, 3, 5, 7 and 10 years has factor exposures of $-1.998$ and $-3.067$ million dollars per basis point of factor score (printed $-1.99$ and $-3.06$). Its two-factor standard deviation is 25.50 (25.40 from the printed exposures; Hull prints 25.45), and $2.326\times25.45=59.2$ is his one-day 99 percent VaR. \emph{Hedge}, original: positions of $-4.80$ million dollars per basis point in the 2-year and $+7.02$ in the 10-year remove both exposures, but factors 3 to 8 leave a standard deviation of 14.88 against 32.30 unhedged across all eight.

\emph{Original, U.S. data.} On month-end Gurkaynak--Sack--Wright zero yields from January~1986 to September~2026 (489 months; \hyperref[data:gsw-treasury-curve]{the Fed's fitted curve}; \citep[\texttt{feds200628.csv}, SVENY01--SVENY30 at 1, 2, 3, 5, 7, 10, 20 and 30 years]{fed-gsw-data}), the first three components of monthly changes explain 85.3, 10.7 and 2.5 percent of variance, with zero, one and two sign changes across maturities. The standard deviation of monthly yield changes is humped: 24.5 basis points at 1 year, 28.7 at 5 years and 25.5 at 30 years. From 2009 to November~2015, near the zero bound, the 1-year yield moved with a standard deviation of only 7.2 basis points against 25.9 at 10 years (27.5 at 1 year in 1986--2007). Multiplying by duration turns this into price volatility rising steeply with maturity, about 0.24 percent a month for a 1-year zero against 7.6 for a 30-year zero.

\qffield{Conditions} Principal components assume symmetric, level-independent shifts and depend on the maturities and period chosen; two factors overstate correlations at the short end, so investors sensitive to specific maturities should monitor key-rate exposures directly (pp.~745--747 and~768--773). Using monthly rather than daily changes, and zero rather than par yields, changes the shares.

\qffield{Pitfalls} Hedging only duration (the first factor) and calling the book immune to curve risk. Assuming the slope factor pivots at the same maturity in every regime. Forming a rate view from the forward curve alone, ignoring term premia (\hyperref[res:term-structure-theories]{theories}).

\qffield{See also} \hyperref[meth:curve-trades-carry-rolldown]{curve trades}; \hyperref[mod:factor-risk-model]{factor risk models}; \hyperref[def:macro-regime-asset-map]{macro regimes}.

\end{qfmethod}
